\clearpage

\section{Clock RG: Derivation and
Normalization}\label{sec-nbcp-clock-rg}

The clock criterion concerns a two-dimensional thermal phase field after
short-distance and nonzero-frequency fluctuations have been integrated
out. It is not an RG equation for the zero-temperature magnon energy.
The assumptions are a stable density-ordered background, a compact phase
\(\phi\equiv\phi+2\pi\), a positive spatial stiffness, weak clock
pinning and dilute integer vortices. Within these assumptions, the
sixfold perturbation can be irrelevant while vortices remain bound, as
in the planar-model analysis of
José et al.~(1977)~\cite{jose1977}. The
derivation below fixes the field, area and fugacity normalizations used
here.

\textbf{Static normalization and constant spatial anisotropy.} Let
\(\mathcal T=k_{\rm B}T\), let \(a_0\) be the matching cutoff of the
phase theory, and take a phase-independent positive symmetric tensor
\(\rho\). The nearest-neighbor distance \(a\), area per spin
\(a_{\rm spin}\) and matching cutoff \(a_0\) are distinct quantities.
The static free energy is

\begin{equation}\phantomsection\label{eq-nbcp-clock-partition}{
F=\int d^2r\left[
\frac12(\nabla\phi)^{\mathsf T}\rho(\nabla\phi)
-\frac{\lambda_p}{a_{\rm spin}}\cos(p\phi-\delta_p)\right],
\qquad Z=\int\mathcal D\phi\,e^{-F/\mathcal T}.
}\end{equation}

Here \(\lambda_p\) is a pinning energy per physical spin at the matching
scale, and \(p\) is an integer harmonic order, not the shorthand for
\(J_{\rm PD}\) used in the preceding SOC calculation. To remove constant
spatial anisotropy, define \(\rho_{\rm eff}=\sqrt{\det\rho}\) and make
the area-preserving change of coordinates

\begin{equation}\phantomsection\label{eq-nbcp-clock-anisotropy-map}{
\mathbf r=L\mathbf R,\qquad
L=\frac{\rho^{1/2}}{(\det\rho)^{1/4}},\qquad
\det L=1,\qquad L^{-1}\rho L^{-\mathsf T}=\rho_{\rm eff}I_2.
}\end{equation}

Then \(d^2r=d^2R\) and the gradient energy is
\(\rho_{\rm eff}\int d^2R\,|\nabla_R\phi|^2/2\). For diagonal \(\rho\),
this means \(x=(\rho_{xx}/\rho_{yy})^{1/4}X\) and
\(y=(\rho_{yy}/\rho_{xx})^{1/4}Y\). The microscopic cutoff shape changes
under this transformation, affecting short-distance matching but not the
Gaussian scaling dimensions. With dimensionless coordinates
\(\mathbf x=\mathbf R/a_0\),

\begin{equation}\phantomsection\label{eq-nbcp-clock-dimensionless-action}{
\mathcal S_0=\frac K2\int d^2x\,(\nabla\phi)^2,
\qquad
\delta\mathcal S=-y_p\int d^2x\,\cos(p\phi-\delta_p),
\qquad
K=\frac{\rho_{\rm eff}}{\mathcal T},\quad
y_p=\frac{a_0^2\lambda_p}{a_{\rm spin}\mathcal T}.
}\end{equation}

\(\mathcal S=F/\mathcal T\) is dimensionless; it is not the spin length
\(S\). For an isotropic tensor \(\rho_{\rm eff}=\rho_s\). A tensor that
depends on \(\phi\) cannot be made constant everywhere by this single
coordinate transformation; that limitation is relevant to SOC Y below.

\textbf{Gaussian correlation function and scaling dimension.} Use
\(\phi(\mathbf x)=\int d^2q\,e^{\mathrm i\mathbf q\cdot\mathbf x}\phi_{\mathbf q}/(2\pi)^2\).
The Gaussian covariance follows by inverting its quadratic kernel:

\begin{equation}\phantomsection\label{eq-nbcp-clock-propagator}{
\mathcal S_0=\frac K2\int\frac{d^2q}{(2\pi)^2}
q^2\phi_{\mathbf q}\phi_{-\mathbf q},\qquad
\langle\phi_{\mathbf q}\phi_{\mathbf q'}\rangle_0
=\frac{(2\pi)^2\delta^{(2)}(\mathbf q+\mathbf q')}{Kq^2}.
}\end{equation}

The zero mode cancels in the phase difference. With
\(R=|\mathbf R|\gg a_0\), angular integration gives

\begin{equation}\phantomsection\label{eq-nbcp-clock-phase-variance}{
\begin{aligned}
W(R)&=\langle[\phi(\mathbf R)-\phi(0)]^2\rangle_0\\
&=\frac{2}{K}\int^{1/a_0}\frac{d^2q}{(2\pi)^2}
\frac{1-\cos(\mathbf q\cdot\mathbf R)}{q^2}\\
&=\frac{1}{\pi K}\int_0^{1/a_0}\frac{dq}{q}[1-J_0(qR)]
=\frac{1}{\pi K}\ln\frac R{a_0}+O(1).
\end{aligned}
}\end{equation}

The last step uses the logarithmic range \(1/R\lesssim q\lesssim1/a_0\);
the cutoff changes the additive constant. Gaussian averaging of a vertex
operator now yields

\begin{equation}\phantomsection\label{eq-nbcp-clock-vertex-dimension}{
\begin{aligned}
\langle e^{\mathrm ip\phi(\mathbf R)}e^{-\mathrm ip\phi(0)}\rangle_0
&=\exp[-p^2W(R)/2]\propto(R/a_0)^{-p^2/(2\pi K)},\\
x_p&=\frac{p^2}{4\pi K},\qquad
\eta=2x_1=\frac{1}{2\pi K}.
\end{aligned}
}\end{equation}

An operator of scaling dimension \(x_p\) has a two-point function
proportional to \(R^{-2x_p}\). The cosine is the sum of charges \(+p\)
and \(-p\) and has this same dimension. Its phase offset \(\delta_p\)
does not affect that dimension.

\textbf{Integrating a momentum shell.} Split \(\phi=\phi_<+\phi_>\),
with the fast modes in \(\Lambda e^{-d\ell}<q<\Lambda\). Their local
variance is

\begin{equation}\phantomsection\label{eq-nbcp-clock-shell-variance}{
\langle\phi_>^2\rangle_0
=\frac1K\int_{\Lambda e^{-d\ell}}^\Lambda
\frac{2\pi q\,dq}{(2\pi)^2q^2}
=\frac{d\ell}{2\pi K}.
}\end{equation}

At first order in \(y_p\), average the perturbation over these modes.
The odd sine average vanishes and the Gaussian cosine average gives

\begin{equation}\phantomsection\label{eq-nbcp-clock-shell-cosine}{
\begin{aligned}
\langle\cos[p(\phi_<+\phi_>)-\delta_p]\rangle_>
&=e^{-p^2\langle\phi_>^2\rangle_0/2}
\cos(p\phi_<-\delta_p)\\
&=e^{-p^2d\ell/(4\pi K)}\cos(p\phi_<-\delta_p).
\end{aligned}
}\end{equation}

Restore the cutoff with \(\mathbf x=e^{d\ell}\mathbf x'\). The measure
contributes \(e^{2d\ell}\), while the Gaussian gradient action is
invariant in two dimensions and the compact angle is not rescaled.
Therefore

\begin{equation}\phantomsection\label{eq-nbcp-clock-pinning-flow}{
y_p'=y_p\exp\left[\left(2-\frac{p^2}{4\pi K}\right)d\ell\right],
\qquad
\frac{dy_p}{d\ell}=\left(2-\frac{p^2}{4\pi K}\right)y_p+\cdots.
}\end{equation}

The two terms have different origins: \(2\) is the area rescaling, and
\(p^2/(4\pi K)\) is the suppression by fast phase fluctuations. Treating
\(K\) as constant gives the exponential solution used in the paper. More
generally, the linearized solution is
\(y_p(\ell)=y_p(0)\exp\{\int_0^\ell[2-p^2/(4\pi K(\ell'))]d\ell'\}\); it
is not a fixed exponential once the stiffness flows.

\textbf{Vortex charge and fugacity.} Compactness permits a winding
\(\oint\nabla\phi\cdot d\mathbf l=2\pi m\), with integer \(m\). Outside
its core, \(\phi=m\arg(X+\mathrm iY)\) and \(|\nabla\phi|=|m|/R\). Hence

\begin{equation}\phantomsection\label{eq-nbcp-clock-vortex-energy}{
E_m=E_{{\rm core},m}
+\frac{\rho_{\rm eff}}2\int_{a_0}^L2\pi R\,dR\,\frac{m^2}{R^2}
=E_{{\rm core},m}+\pi\rho_{\rm eff}m^2\ln(L/a_0).
}\end{equation}

A vortex position contributes a measure \(d^2R/a_0^2\), or entropy
\(2k_{\rm B}\ln(L/a_0)\) for one defect. Its dimensionless weight
therefore scales as \(y_m(L/a_0)^{2-\pi Km^2}\), where \(y_m\) includes
the core weight and its local prefactor. This identifies
\(x_{v,m}=\pi Km^2\). The isolated-vortex argument extracts the scaling
dimension; a periodic system has net charge zero and contains neutral
vortex configurations. For charges \(m=\pm1\) with fugacity \(y_v\) per
charge,

\begin{equation}\phantomsection\label{eq-nbcp-clock-vortex-flow}{
\frac{dy_v}{d\ell}=(2-\pi K)y_v+\cdots.
}\end{equation}

Vortices are irrelevant at the Gaussian line only if \(K>2/\pi\). A
smooth-wave LSWT calculation cannot determine \(y_v\): the vortex core
involves configurations outside a small-amplitude expansion about one
uniform state.

\textbf{Why the stiffness also flows.} The preceding eigenvalues are
exact at the Gaussian line, but keeping \(K\) fixed discards the leading
feedback of defects. To display its normalization, use normal-ordered
vertices at the cutoff and a real-space dipole shell
\(a_0<r<a_0e^{d\ell}\). Each pinning charge in
\(\exp[y_p\int\cos(p\phi)]\) has fugacity \(y_p/2\). The neutral part of
the second-order cumulant contributes

\begin{equation}\phantomsection\label{eq-nbcp-clock-neutral-dipole}{
\delta\mathcal S_{\rm neutral}
=-\frac{y_p^2}{4a_0^4}\int d^2R\int_{\rm shell}d^2r
\left(\frac{a_0}{r}\right)^{2x_p}
\cos\!\left[p\left(\phi(\mathbf R+\mathbf r/2)-\phi(\mathbf R-\mathbf r/2)\right)\right].
}\end{equation}

The constant term changes the free energy. Expanding the phase
difference as \(\mathbf r\cdot\nabla\phi\) and the cosine as
\(1-p^2(\mathbf r\cdot\nabla\phi)^2/2\) gives

\begin{equation}\phantomsection\label{eq-nbcp-clock-electric-screening}{
\begin{aligned}
\int_{\rm shell}d^2r\,r_\mu r_\nu
\left(\frac{a_0}{r}\right)^{2x_p}
&=\pi a_0^4\delta_{\mu\nu}\,d\ell+O(d\ell^2),\\
\delta\mathcal S_{\rm grad}
&=\frac{\pi p^2y_p^2}{8}d\ell\int d^2R\,|\nabla\phi|^2,
\qquad
\left.\frac{dK}{d\ell}\right|_p=\frac{\pi p^2}{4}y_p^2.
\end{aligned}
}\end{equation}

Pinning increases \(K\) at this order. The vortex contribution can be
obtained from the dual Gaussian field \(\theta\). The neutral vortex gas
has pair weight \(\prod_{i<j}(r_{ij}/a_0)^{2\pi K m_i m_j}\), reproduced
by the dual action

\begin{equation}\phantomsection\label{eq-nbcp-clock-dual-action}{
\mathcal S_{\rm dual}
=\frac{K_d}{2}\int d^2R\,|\nabla\theta|^2
-\frac{2y_v}{a_0^2}\int d^2R\,\cos\theta,
\qquad K_d=\frac{1}{4\pi^2K}.
}\end{equation}

Indeed, \(x_{e^{\mathrm i\theta}}=1/(4\pi K_d)=\pi K\). Applying the
preceding dipole result with \(p=1\) and cosine amplitude \(2y_v\) gives
\(dK_d/d\ell=\pi y_v^2\), or \(dK^{-1}/d\ell=4\pi^3y_v^2\). The dual
field is an alternative representation of vortex effects, not an
additional independent phase field. Combining the leading terms yields,
in this dipole-fugacity convention,

\begin{equation}\phantomsection\label{eq-nbcp-clock-coupled-flow}{
\begin{aligned}
\frac{dy_p}{d\ell}&=\left(2-\frac{p^2}{4\pi K}\right)y_p+\cdots,\\
\frac{dy_v}{d\ell}&=(2-\pi K)y_v+\cdots,\\
\frac{dK^{-1}}{d\ell}&=4\pi^3y_v^2-
\frac{\pi p^2}{4K^2}y_p^2+\cdots.
\end{aligned}
}\end{equation}

The quadratic coefficients require the stated fugacity normalization and
cutoff convention; rescaling either fugacity changes them. The linear
scaling dimensions and the opposing effects of vortex screening and
pinning do not depend on that choice. Higher orders generate additional
harmonics and mixed operators. The displayed equations are a
dilute-defect, weak-pinning truncation, not a quantitative flow for
arbitrary SOC.

The algebraic-window inequalities and their physical conditions are
stated in Section~\ref{sec-nbcp-clock}. They use the stiffness at the
relevant scale and must not be converted to transition temperatures by
inserting an unmatched classical coefficient.
